\chapter{Experimental Validation and Results}
\chaplabel{ch:validation}
\label{sec:testing}

\section*{Introduction}
The preceding chapters described the objectives set for the work, the architecture designed to meet them, and the implementation that realized it. This chapter establishes whether the result satisfies those objectives and the requirements derived from them. It states the strategy by which the engine was tested, the bench on which the tests were run, and the means by which the reconstructed trace was checked for correctness. It then reports the defects that appeared only on hardware, the results obtained against each requirement, and the limitations that remain together with the directions they suggest.


\section{Test Strategy}
Testing was organized in three tiers, distinguished by what each requires.

The first tier requires nothing beyond the built engine. It covers the protocol handshake, the cancellation path, and the components that operate on stored data: the device description parser, the static program analysis, the trace decoder and the reconstruction transformations. This tier runs during ordinary development.

The second tier requires a device. It drives the engine at the protocol level with a script that sends DAP requests and checks responses, covering launching, attaching, stepping, breakpoints, peripheral access and trace.

The third tier requires a device and the frontend. It drives the extension inside VS Code and checks the behavior a developer would observe, including that the views populate and that state is refreshed after a reset. It is exercised by hand.

No unit testing framework is used. Tests in the first tier are plain executables that return a status, keeping the build free of an additional dependency.

Concurrency is covered separately. A test publishes on the event bus from several threads at once, so that a defect in its locking appears as a detected race. The engine was additionally built with a thread sanitizer and driven against the device, which is how the lock-order inversion of Section \ref{subsec:lldb_impl} was identified. Table \ref{tab:tests} in the annexes lists the tests.


\section{Hardware Test Bench}
\label{subsec:bench}
The bench consists of a development board carrying the target device together with an on-board probe, connected to the host by a single cable. It contains no trace probe and no connection to trace pins. Every trace result reported below was obtained over the same connection used for ordinary debugging, which is the objective stated in Chapter \ref{ch:specification}.

A small firmware project serves as the program under examination, built with debug information and with the optimization settings a developer would ordinarily use.

Every result reported in this chapter was obtained in a session that the engine launched, since trace, live memory observation and peripheral inspection all depend on the direct path.

The pipeline was additionally exercised on an STM32H7 and an STM32N6, the second of which implements a different trace unit on a later processor. The base addresses and the OpenOCD configuration script changed. Neither driver was modified, since each reads the capabilities of the component it drives from that component's own identification registers.


\section{Trace Correctness Validation}
\label{subsec:trace_validation}
Correctness was checked against a stored capture, against the device, and against commercial debuggers.

The decoding and reconstruction stages are checked against a stored capture. The capture, the program image and the trace unit register values recorded from the device are held as test data. The test decodes the capture and asserts that the number of executed instructions equals a fixed expected value, that every reconstructed address corresponds to an instruction in the static analysis, and that at least one function block carries a resolved function name. Holding the register values as data allows the test to reproduce the configuration the hardware actually had.

The acquisition stages are checked against the device. A test asserts that the trace components are found, that the register values read from the trace unit match those recorded, and that a capture delivers exactly two callbacks per halt: the synchronization barrier followed by the payload.

The reconstruction was compared against the output of three commercial debuggers, each driving the probe of its own vendor. The decoded listings were compared instruction by instruction, and the sequences agreed. The differences were confined to presentation.

A capture obtained by a probe of another vendor was decoded by the same pipeline. Once the synchronization framing is discounted, the byte stream was identical to the one drained from the on-chip sink.


\section{Defects Found on Hardware}
\label{subsec:field_issues}
Driving the result against a real device produced defects that no test on stored data would have found.

\begin{table}[htbp]
    \centering
    \caption{Defects found on hardware}
    \label{tab:defects}
    \begin{tabularx}{\linewidth}{| @{}>{\hsize=0.85\hsize\bfseries}Y | >{\hsize=1.65\hsize}Y | >{\hsize=1.1\hsize}Y | >{\hsize=0.4\hsize}Y@{} |}
        \hline
        Symptom & \textbf{Cause} & \textbf{Fix} & \textbf{Section} \\
        \hline
        A breakpoint on a source line resolved into a vendor header & Several source locations overlap at an inlined call site and the innermost was selected & Walk outward to the outermost inlined block and use the call site recorded there & \ref{subsec:breakpoints_impl} \\
        \hline
        Breakpoints could not be set in assembly files & The frontend had not declared that it supports breakpoints in that language & Declare the language & \ref{subsec:breakpoints_impl} \\
        \hline
        A restart left a stale view of the program & The reset is performed outside the conventional path, so LLDB was never told & The resynchronization mechanism & \ref{subsec:reset_impl} \\
        \hline
        Variables and scopes were stale after a restart & The same cause & The same mechanism & \ref{subsec:reset_impl} \\
        \hline
        Trace was inconsistent and often empty & The components were not armed on every transition that resumes the processor, and reset was mishandled & Couple both drivers to the run state of the target & \ref{subsec:drivers} \\
        \hline
        Stepping did not halt where expected & The internal breakpoints LLDB places for a step were software breakpoints in flash, which never trigger & Force those breakpoints to use comparators & \ref{subsec:breakpoints_impl} \\
        \hline
    \end{tabularx}
\end{table}


\section{Results}
The objective stated in Chapter \ref{ch:specification} was met. Instruction trace is captured from an on-chip sink, extracted over the ordinary debug connection, decoded, reconstructed into a sequence of executed instructions and attributed to source locations, and presented inside a development tool during a live debug session. No trace pins and no specialized probe are used.

The debugger that carries this capability provides breakpoints with automatic selection between software and hardware mechanisms, stepping at source and instruction granularity, variable and expression inspection, memory reading and writing, disassembly, device reset, and the extensions covering trace, live memory observation and peripheral inspection. The engine builds and runs on two host platforms.

\begin{table}[htbp]
    \centering
    \caption{Verification of the requirements}
    \label{tab:verification}
    \begin{tabularx}{\linewidth}{| @{}>{\hsize=0.35\hsize\bfseries}Y | >{\hsize=1.55\hsize}Y | >{\hsize=1.1\hsize}Y@{} |}
        \hline
        ID & \textbf{How it was established} & \textbf{Status} \\
        \hline
        FR1 & Trace source and sink located and their registers read back against the device & Met \\
        \hline
        FR2 & Capture delivery observed on every halt, two callbacks per halt & Met \\
        \hline
        FR3 & Every result in this chapter obtained over the debug connection alone & Met \\
        \hline
        FR4 & Decode of the stored capture, Figure \ref{fig:bytes_to_elements} & Met \\
        \hline
        FR5 & Every reconstructed address resolves to an instruction in the image, the counts match the frozen baseline of Table \ref{tab:measured}, and the reconstructed sequence was compared against three vendor debuggers, Section \ref{subsec:trace_validation} & Met \\
        \hline
        FR6 & Function blocks carrying resolved names in the same test & Met \\
        \hline
        FR7 & Gaps counted and represented explicitly in the same test & Met \\
        \hline
        FR8 & Trace views populated during a live session, observed in the development tool & Met \\
        \hline
        FR9 & \texttt{\scriptsize stepping\_smoke\_test} against the device, covering stepping over, into and out of a call, and breakpoints & Met \\
        \hline
        FR10 & Memory reading exercised by \texttt{\scriptsize launch\_smoke\_test} against the device, and register access exercised on both paths & Met \\
        \hline
        FR11 & \texttt{\scriptsize peripheral\_smoke\_test} against the device, covering enumeration, reading and writing & Met \\
        \hline
        FR12 & Disassembly, scopes and variables exercised by \texttt{\scriptsize launch\_smoke\_test} against the device & Met \\
        \hline
        FR13 & Reset followed by the resynchronization of Section \ref{subsec:reset_impl}, observed against the device & Met \\
        \hline
        NFR1 & The bench carries no trace probe and no connection to trace pins, Section \ref{subsec:bench} & Met \\
        \hline
        NFR2 & The engine builds and runs on both host platforms & Met \\
        \hline
        NFR3 & The frontend is an extension speaking the standard protocol & Met \\
        \hline
    \end{tabularx}
\end{table}

\begin{table}[htbp]
    \centering
    \caption{Measured figures for one capture}
    \label{tab:measured}
    \begin{tabularx}{\linewidth}{| @{}>{\hsize=1.15\hsize\bfseries}Y | >{\hsize=0.85\hsize}Y@{} |}
        \hline
        Quantity & \textbf{Value} \\
        \hline
        Interval traced & Reset handler to the entry of \texttt{main} \\
        \hline
        Trace produced by the device & 128 bytes, eight formatter frames \\
        \hline
        Host-inserted synchronization barrier & 16 bytes, one frame \\
        \hline
        Bytes presented to the decoder & 144 bytes, nine formatter frames \\
        \hline
        Elements retained & 35 \\
        \hline
        Executed instructions reconstructed & 114 \\
        \hline
        Function blocks & 6 \\
        \hline
        Start-of-capture markers & 2 \\
        \hline
        Discontinuities within the traced execution & 0 \\
        \hline
        Executed instructions per byte produced by the device & 0.89 \\
        \hline
    \end{tabularx}
\end{table}

The barrier is emitted by the extraction driver ahead of every payload, so it never occupied the sink and is counted apart from the trace the device produced. Both figures are given because the first states what the buffer held and the second states what the decoder consumed.

The last row answers what a buffer represents. At the density this capture exhibits, the two kilobytes of the sink hold on the order of eighteen hundred executed instructions. The figure is indicative, since the density depends on how often the program branches.


\section{Limitations and Future Work}
\label{subsec:limits}
\label{subsec:trace_limits}

\begin{table}[htbp]
    \centering
    \caption{Limitations of the trace pipeline}
    \label{tab:trace_limits_tab}
    \begin{tabularx}{\linewidth}{| @{}>{\hsize=0.75\hsize\bfseries}Y | >{\hsize=1.25\hsize}Y@{} |}
        \hline
        Limitation & \textbf{What it means} \\
        \hline
        Component discovery & The trace source and the trace sink are located by name among the objects created by configuration, and their base addresses are stated there. Supporting a new device requires a configuration file to be written. \\
        \hline
        One protocol version & The drivers and the decoder handle the version implemented by the devices used. No device reachable by the on-chip route is lost, since every STM32 device carrying an on-chip trace sink implements that version, as Table \ref{tab:stm32_debug} states. \\
        \hline
        One source and one sink & The first of each is selected, so a device offering several is not fully exploited. Tracing several cores is outside the scope of a proof of concept on a single-core target. \\
        \hline
        Timing discarded & Timestamps and cycle counts are decoded and dropped at the sink, so no timing analysis is possible without changing that component. \\
        \hline
        Exceptions not projected & The elements are retained and no consumer presents them, so an exception appears in the record as a change of address. \\
        \hline
        No trace filtering & The implementation traces without filtering, so the volume of trace is bounded by the size of the buffer alone. On these devices filtering would be driven from the processor comparators. \\
        \hline
        Return classification & LLVM's analysis facility does not recognize the instruction forms this architecture emits for a return, so those forms are matched on the printed instruction. \\
        \hline
        Offline decoding not exposed & The stages require only a capture, a program image and the recorded register values. The standalone entry point that made this available during development was not carried forward. \\
        \hline
    \end{tabularx}
\end{table}

Further limitations concern the engine. The two disassembly implementations have not been unified. Several requests drive LLDB outside the serialization of Section \ref{subsec:lldb_impl}. Capabilities depending on the direct path are unavailable when attaching to an existing session. The frontend is not distributed.

Several directions extend the work. Retaining the timing elements currently discarded would enable profiling from the same captures. Presenting exceptions explicitly would complete the reconstruction of execution across interrupt entry. Exposing the filtering capability of the trace unit would allow the buffer to cover a chosen region. Reading the component topology from the ROM tables would remove the per-device configuration file.


\section{Conclusion}
The result was measured against the objectives stated in Chapter \ref{ch:specification} and against the requirements derived from them. Correctness was confirmed against stored captures, against the device, and against commercial debuggers, whose reconstructed listings agreed instruction by instruction. Every functional and non-functional requirement was met, and instruction trace was obtained over the ordinary debug connection with no trace pins and no probe beyond the one embedded on the board. Given the demonstrated approach, the limitations open directions for further work.
